\scanpage{62}
\uline{Remarque} Supposons $k$ principal. Alors la catégorie $P\simeq P_{\bullet}$ \deleted{\unsure{est ab.}} admet des \deleted{$\Sigma$} noyaux (et des Conoyaux, \deleted{\unsure{?|$\Sigma$}} mais un bimorphisme n'est pas néc. un iso). \deleted{\unsure{On fera attention que les mor.}} Il est \unsure{m.|il} vrai que si $N$ est le noyau de $A\longrightarrow B$, alors $N\longrightarrow A$ est un $\Sigma$-mono, i.e. s'insère dans une suite $\Sigma$-exacte $0\longrightarrow N\longrightarrow A\longrightarrow A'\longrightarrow0$, mais le mor. $A'\longrightarrow B$ ne sera pas en général un $\Sigma$-mono. \deleted{Aussi} On ne peut interpréter $(P,\deleted{\unsure{\Psi^0}}\mathbb{C}^{\bullet},\deleted{\unsure{\Psi^0,\varepsilon}}\inserted{\Sigma}\deleted{\unsure{,\Sigma}})$ comme sol. d'un pb universel en appliquant tels quels les développements précédents: \deleted{Mais \deleted{\unsure{disons}} \inserted{\unsure{on peut}} définir de façon naturelle, pour une $B$ $k$-linéaire, la notion d'objet \uline{plat}, et alors} le syst. ci-dessus est universel pour les $B$ \inserted{$k$-additifs} munis d'un complexe \deleted{de cochaînes \inserted{\unsure{libres}} plat, \deleted{\unsure{des}} d'un. \deleted{sys} \inserted{classe} $\Sigma$ de suites courtes, telles que tt hom $u\inserted{{}:A\longrightarrow B}$ de $B$ ait un noyau tel que $\operatorname{Ker}u\longrightarrow A$ soit $\in\Sigma$} $C^{\bullet}$ et d'une classe $\Sigma_C$ de suites exactes, telle que tt morphisme $u:A\longrightarrow B$ de \unsure{$\mathcal B$|$B$} admette un noyau $K$ tel que $K\longrightarrow A$ soit un $\Sigma_C$-mono: (on prend comme morphismes entre tels syst. ceux qui transforment $C^{\bullet}$ en $C'^{\bullet}$ et qui sont $\Sigma_{\unsure{i|1}}$-exacts à \uuline{g}).
% archive sheet number: 61
